\section*{Capítulo \ref{ch:g9:elements-universe-earth} --- Os elementos no Universo e na Terra}

\begin{solution}{exo:g9:elements-universe-earth:1}
O hidrogênio (cerca de \qty{74}{\percent}) e o hélio (cerca de
\qty{25}{\percent}).
\end{solution}

\begin{solution}{exo:g9:elements-universe-earth:2}
O oxigênio, nos dois casos.
\end{solution}

\begin{solution}{exo:g9:elements-universe-earth:3}
Dentro das estrelas, e nas explosões delas.
\end{solution}

\begin{solution}{exo:g9:elements-universe-earth:4}
Cerca de \qty{3.5}{\percent}.
\end{solution}

\begin{solution}{exo:g9:elements-universe-earth:5}
$46.6 + 27.7 = \qty{74.3}{\percent}$: cerca de três quartos.
\end{solution}

\begin{solution}{exo:g9:elements-universe-earth:6}
Eles são gases muito leves: escaparam da Terra jovem para o espaço (e o
hélio não forma compostos que pudessem tê-lo retido nas rochas).
\end{solution}

\begin{solution}{exo:g9:elements-universe-earth:7}
$0.23 \times 70 = \qty{16.1}{kg}$.
\end{solution}

\begin{solution}{exo:g9:elements-universe-earth:8}
Os \omterm{def:g7:atoms-and-molecules:atom}{átomos} de oxigênio são 16 vezes mais pesados que os de hidrogênio:
menos \omterm{def:g7:atoms-and-molecules:atom}{átomos} de oxigênio pesam mesmo assim mais.
\end{solution}

\begin{solution}{exo:g9:elements-universe-earth:9}
Cerca de $0.081 \times 1000 = \qty{81}{kg}$ de alumínio.
\end{solution}

\begin{solution}{exo:g9:elements-universe-earth:10}
O oxigênio e o cálcio (e, além dos oito principais elementos da crosta,
o hidrogênio, o carbono e o fósforo estão presentes na crosta em
pequenas quantidades).
\end{solution}

\begin{solution}{exo:g9:elements-universe-earth:11}
$40 \times 1.67 \times 10^{-27} = \qty{6.68e-26}{kg}$;
$1.06 / 6.68 \times 10^{-26} \approx 1.6 \times 10^{25}$ \omterm{def:g7:atoms-and-molecules:atom}{átomos} de cálcio.
\end{solution}

\begin{solution}{exo:g9:elements-universe-earth:12}
Os \omterm{def:g7:atoms-and-molecules:atom}{átomos} não são criados nem destruídos pelas \omterm{def:g5:irreversible-changes:chemical-change}{transformações químicas}:
eles só passam de um composto para outro (rocha, gás, açúcar, corpo
vivo). Os \omterm{def:g7:atoms-and-molecules:atom}{átomos} de carbono da Terra são usados muitas e muitas vezes.
\end{solution}

\begin{solution}{pb:g9:elements-universe-earth:1}
\textbf{1.} $0.61 \times 70 = \qty{42.7}{kg}$.

\textbf{2.} Carbono $0.23 \times 70 = \qty{16.1}{kg}$; hidrogênio
$0.10 \times 70 = \qty{7.0}{kg}$.

\textbf{3.} $61 + 23 + 10 = \qty{94}{\percent}$.

\textbf{4.} A água, \ce{H2O}.

\textbf{5.} $16 \times 1.67 \times 10^{-27} = \qty{2.67e-26}{kg}$.

\textbf{6.} Carbono: $12 \times 1.67 \times 10^{-27} = \qty{2.00e-26}{kg}$;
hidrogênio: \qty{1.67e-27}{kg}.

\textbf{7.} 16.

\textbf{8.} Um \omterm{def:g9:inside-the-atom:nucleus}{elétron} é cerca de 1836 vezes mais leve que um \omterm{def:g9:inside-the-atom:proton}{núcleon}:
os \omterm{def:g9:inside-the-atom:nucleus}{elétrons} representam menos de um milésimo da massa.

\textbf{9.} $7.0 / 1.67 \times 10^{-27} \approx 4.2 \times 10^{27}$
\omterm{def:g7:atoms-and-molecules:atom}{átomos} de hidrogênio.

\textbf{10.} $16.1 / 2.00 \times 10^{-26} \approx 8.0 \times 10^{26}$
\omterm{def:g7:atoms-and-molecules:atom}{átomos} de carbono.

\textbf{11.} $42.7 / 2.67 \times 10^{-26} \approx 1.60 \times 10^{27}$
\omterm{def:g7:atoms-and-molecules:atom}{átomos} de oxigênio.

\textbf{12.} A nossa contagem, $1.60 \times 10^{27}$, concorda com a
estimativa $1.61 \times 10^{27}$: um corpo de \qty{70}{kg} contém cerca
de $1.6 \times
10^{27}$ \omterm{def:g7:atoms-and-molecules:atom}{átomos} de oxigênio.
\end{solution}
